\chapter{Theoretical Background}
\label{chap:background}

This chapter presents the theoretical foundations underpinning the Pseudo++ toolchain.
It covers the Monaco editor and its Monarch tokenisation model, the mechanics of
incremental parsing via Tree-sitter, semantic analysis and static checking, the
distinction between dynamic and static typing and its implications for transpilation,
the structured-programming basis for loop-equivalence transformations, the WebAssembly
execution model, and the step-through debugging model.

% ─────────────────────────────────────────────────────────────────────────────
\section{Monaco Editor and Monarch Tokenisation}
\label{sec:monaco-monarch}

\subsection{Monaco Editor}

Monaco Editor \cite{monaco} is the open-source code-editing component that powers
Visual Studio Code.  It is written in TypeScript and runs entirely in the browser,
providing a full-featured editing experience --- line numbers, bracket matching, glyph
margins, decorations, and theming --- without a server-side component.

Monaco separates two language-related concerns:

\begin{itemize}
  \item \textbf{Tokenisation} --- assigning a visual token class (keyword, string,
        operator, comment, \ldots) to each contiguous run of characters for the
        purpose of syntax highlighting.
  \item \textbf{Language services} --- richer, semantic operations such as
        auto-completion, hover information, and diagnostics, which require a deeper
        understanding of the language structure.
\end{itemize}

Pseudo++ uses Monaco for its web editor, registering a custom language
(\texttt{pseudocode}) and providing both a Monarch tokeniser (Section~\ref{subsec:monarch})
and a language configuration (bracket pairs, auto-closing, comment tokens).

\subsection{The Monarch Tokeniser Model}
\label{subsec:monarch}

Monarch \cite{monarch} is Monaco's declarative tokeniser specification language.
A Monarch grammar is a JavaScript object describing a \textit{finite-state transducer}:
it has a set of \textit{states}, and each state contains an ordered list of
\textit{rules} of the form $(\mathit{regex}, \mathit{action})$.

The tokeniser processes the input line by line.  In state $q$, it tries each rule's
regex against the current position; the first match fires the associated action, which
emits a token class and optionally transitions to a new state.  This gives Monarch the
expressive power of a \textit{regular language} per state, with the ability to
represent richer patterns (e.g.\ multi-line strings) through state transitions.

For the Pseudo++ language, the single \texttt{root} state suffices.  Rules recognise,
in priority order: comments (\texttt{\#\ldots}), string literals, numbers, the
swap and assignment operators (\texttt{<->}, \texttt{<-}), comparison operators, the
floor brackets (\texttt{[}, \texttt{]}), and finally identifiers --- which are
dispatched through case-insensitive keyword tables
(\texttt{controlKeywords}, \texttt{flowKeywords}, \texttt{builtinFunctions},
\texttt{logicalOperators}) to assign the appropriate token class, falling back to
\texttt{variable} for unknown identifiers.

\subsection{Monarch vs.\ Tree-sitter}

Monarch and Tree-sitter serve complementary but distinct roles in Pseudo++:

\begin{itemize}
  \item \textbf{Monarch} runs entirely in JavaScript inside Monaco, is
        regex-based, and operates line by line.  It is fast enough to re-tokenise
        the visible viewport on every keystroke, producing token classes consumed
        directly by Monaco's renderer.  It has no notion of nesting, scope, or
        grammatical structure.

  \item \textbf{Tree-sitter} runs as a WebAssembly module compiled from C, performs
        full incremental GLR parsing, and produces a concrete syntax tree (CST)
        encoding the complete grammatical structure of the program.  It is used
        exclusively by the interpreter and transpiler for execution and code
        generation.
\end{itemize}

The two tools are never in conflict: Monarch makes the editor \textit{look}
right; Tree-sitter makes execution \textit{work} right.

% ─────────────────────────────────────────────────────────────────────────────
\section{Syntactic Analysis and Incremental Parsing}
\label{sec:parsing}

\subsection{Classical Parsing Approaches}

Given a grammar $G$ and an input string $w$, a \textit{parser} determines whether
$w \in L(G)$ and, if so, constructs a parse tree witnessing the derivation.  Two
broad families of algorithms exist:

\begin{itemize}
  \item \textbf{Top-down parsers} (e.g.\ LL($k$), recursive descent) expand the
        start symbol top-down.  They are simple to implement and produce intuitive
        error messages, but require the grammar to be left-recursion-free and
        left-factorable.

  \item \textbf{Bottom-up parsers} (e.g.\ LR($k$), LALR(1)) shift tokens onto a
        stack and reduce them by grammar rules.  They handle a strictly larger class
        of grammars than LL parsers and are the basis for most production parser
        generators (\texttt{yacc}, \texttt{bison}, \texttt{menhir}).
\end{itemize}

\subsection{Tree-sitter and Incremental GLR Parsing}

Tree-sitter \cite{treesitter} is an open-source parsing framework that generates
Generalised LR (GLR) parsers from a grammar described in a JavaScript DSL.  Its
defining feature is \textit{incremental re-parsing}: when the source text changes,
the parser does not restart from scratch but instead re-uses unchanged subtrees from
the previous parse tree.

Formally, let $T$ be a concrete syntax tree (CST) for source $s$, and let $s'$ be
the result of applying an edit $\delta$ to $s$.  Tree-sitter maintains a set of
\textit{valid subtrees} --- nodes whose byte ranges are unaffected by $\delta$ ---
and constructs $T'$ by re-parsing only the minimal affected region, grafting in the
cached subtrees.  The amortised complexity of an edit of size $k$ is $O(k \log n)$
where $n = |s|$ \cite{treesitter-thesis}.  In Pseudo++, Tree-sitter is invoked
each time the user requests execution or transpilation, giving an always up-to-date
CST without redundant full re-parses.

\subsection{Concrete vs.\ Abstract Syntax Trees}

Tree-sitter produces \textit{concrete syntax trees} (CSTs), which contain every
token including punctuation and delimiter tokens.  By contrast, an
\textit{abstract syntax tree} (AST) retains only semantically meaningful nodes.

Pseudo++ works directly on the CST exposed by Tree-sitter's C API.  Node types
are string-tagged (\texttt{ts\_node\_type}), and named fields (e.g.\
\texttt{condition}, \texttt{var}, \texttt{start}) allow the interpreter and
transpiler to locate child nodes by role rather than by index.  This design avoids
a separate lowering pass to an AST while retaining the benefits of a typed node API.

\subsection{Error Recovery}

Tree-sitter employs a built-in error-recovery strategy: when the parser encounters
an unexpected token, it inserts \texttt{ERROR} and \texttt{MISSING} nodes into the
tree rather than aborting.  The resulting tree is always complete (every byte of
input appears in some node) and serves as a basis for incremental syntax-error
diagnostics.  Pseudo++'s parser module traverses the tree looking for these sentinel
nodes and reports their location to the user.

% ─────────────────────────────────────────────────────────────────────────────
\section{Semantic Analysis and Static Checking}
\label{sec:semantic-analysis}

\subsection{Beyond Syntax}

Syntactic analysis determines whether a program is grammatically well-formed;
semantic analysis determines whether it is \textit{meaningful}.  A program can
be syntactically valid yet semantically erroneous --- for example, using a
variable before it is assigned, or passing a string where a number is expected.

Semantic analysis operates on the parse tree produced by the parser and
enforces constraints that cannot be expressed in a context-free grammar.  It
is typically organised into two phases:

\begin{itemize}
  \item \textbf{Scope analysis} --- resolves identifiers to their declarations
        and detects uses of undeclared or uninitialised variables.  This
        requires maintaining a \textit{symbol table}: a mapping from identifier
        names to their declaration context and inferred type.

  \item \textbf{Type checking} --- verifies that operations are applied to
        operands of compatible types (e.g.\ that arithmetic operators receive
        numeric operands) and that assignments are type-consistent.
\end{itemize}

\subsection{Linting}

A \textit{linter} is a tool that performs static semantic analysis without
executing the program, reporting potential errors as diagnostics.  The term
originates with the Unix \texttt{lint} utility for C, which flagged type
mismatches and undefined behaviour that the compiler accepted silently.

Modern linters are oriented towards \textit{warnings} rather than hard errors:
they flag patterns that are likely mistakes --- an uninitialised variable, a
condition that is always true, unreachable code --- without rejecting the
program outright.  This is well suited to an educational context, where partial
or exploratory programs are common and immediate, non-fatal feedback is more
useful than a hard compile failure.

% ─────────────────────────────────────────────────────────────────────────────
\section{Interpretation and Transpilation}
\label{sec:interpretation-transpilation}

\subsection{Interpreters}

An \textit{interpreter} is a program that directly executes a source program without
producing standalone machine code.  Interpreters can be categorised by the
intermediate representation they operate on:

\begin{enumerate}
  \item \textbf{Tree-walking interpreters} traverse the AST/CST and evaluate each
        node in turn.  They are simple to implement and straightforward to extend, at
        the cost of performance compared to bytecode approaches.

  \item \textbf{Bytecode interpreters} first compile to an intermediate bytecode
        (e.g.\ CPython's \texttt{.pyc}), then execute it in a virtual machine.
        This separation allows optimisations such as constant folding.

  \item \textbf{JIT compilers} (e.g.\ LuaJIT, V8) profile execution and
        dynamically compile hot paths to native machine code.
\end{enumerate}

Pseudo++ uses a \textit{stack-based tree-walking interpreter}: an explicit execution
stack mirrors the nesting of syntactic constructs, converting what would otherwise be
recursive C calls into an iterative state machine.  Each stack frame carries a
\textit{phase} counter that encodes progress through the frame's sub-steps, making
the execution granularity uniform --- exactly one statement per step --- which is a
prerequisite for the debugger's step-by-step inspection (Section~\ref{sec:debugger}).

\subsection{Transpilation}

\textit{Transpilation} (source-to-source compilation) is a form of compilation where
the target is another high-level language.  Well-known examples include
TypeScript-to-JavaScript (\texttt{tsc}), Sass-to-CSS, and Babel's ES2022-to-ES5
lowering.  In the educational domain, transpilation from a didactic notation to a
production language serves as a scaffolding device: students write in familiar
pseudocode and observe the idiomatic equivalent in C or Pascal.

Pseudo++'s transpiler performs a single-pass tree walk over the CST, emitting
target-language code for each node.  A pre-pass (\texttt{transpiler\_collect}) scans
all assignment targets to produce the forward variable declarations required by C
and Pascal, since both languages mandate that all variables are declared before use.

\subsection{Dynamic and Static Typing}

Programming languages differ in \textit{when} type information is checked.
\textit{Statically typed} languages (C, C++, Pascal) require all variables to
be declared with an explicit type before use; the compiler verifies type
correctness at compile time and rejects ill-typed programs.
\textit{Dynamically typed} languages attach type information to values at
runtime; a variable may hold an integer in one iteration and a string in the
next.

The Romanian baccalaureate pseudocode is dynamically typed by convention:
variables are not declared, and their types are inferred from the values
assigned to them.  Transpiling to C or Pascal therefore requires a
\textit{type inference} pre-pass that analyses all assignment targets and
deduces a static type for each variable before code emission begins.

\subsection{The Pseudo++ Processing Pipeline}
\label{subsec:pipeline}

After parsing, the CST is passed to one of two independent backends depending on
the user's intent.  Figure~\ref{fig:pipeline} shows the complete pipeline.

\begin{figure}[h]
\centering
\resizebox{\linewidth}{!}{%
\begin{tikzpicture}[
  node distance  = 7mm and 22mm,
  box/.style     = {draw, rounded corners=3pt, minimum width=34mm,
                    minimum height=9mm, align=center, font=\small},
  arr/.style     = {-{Stealth[length=4pt]}, thick},
  lbl/.style     = {font=\footnotesize\itshape},
]
  \node[box] (src)   {Canonical Source};
  \node[box, below=of src]   (fmt)   {Formatter};
  \node[box, below=of fmt]   (ascii) {Normalised Source (ASCII)};
  \node[box, below=of ascii] (parse) {Parser (Tree-sitter)};

  \node[box, below left =14mm and 72mm of parse] (interp) {Interpreter};
  \node[box, below left =14mm and 24mm of parse] (linter) {Check};
  \node[box, below right=14mm and 24mm of parse] (debug)  {Debugger};
  \node[box, below right=14mm and 72mm of parse] (trans)  {Transpiler};

  \node[box, below=of interp] (out1) {Program Output};
  \node[box, below=of linter] (out2) {Diagnostics};
  \node[box, below=of debug]  (out3) {Step State};
  \node[box, below=of trans]  (out4) {C / C++ / Pascal};

  \draw[arr] (src)   -- (fmt);
  \draw[arr] (fmt)   -- (ascii);
  \draw[arr] (ascii) -- (parse);
  \draw[arr] (parse) -- node[lbl, above left  ]{run}       (interp);
  \draw[arr] (parse) -- node[lbl, left        ]{check}     (linter);
  \draw[arr] (parse) -- node[lbl, right       ]{debug}     (debug);
  \draw[arr] (parse) -- node[lbl, above right ]{transpile} (trans);
  \draw[arr] (interp) -- (out1);
  \draw[arr] (linter) -- (out2);
  \draw[arr] (debug)  -- (out3);
  \draw[arr] (trans)  -- (out4);
\end{tikzpicture}%
}
\caption{The Pseudo++ processing pipeline.  After normalisation and parsing, the
CST is handed to the interpreter, static analyser, debugger, or transpiler depending on the
requested operation.}
\label{fig:pipeline}
\end{figure}

% ─────────────────────────────────────────────────────────────────────────────
\section{Equivalence of Repetitive Control Structures}
\label{sec:loop-equivalence}

\subsection{Structured Programming and the Böhm-Jacopini Theorem}

The theoretical basis for loop equivalence lies in \textit{structured programming
theory}.  Böhm and Jacopini proved in 1966 \cite{bohm-jacopini} that any computable
function can be expressed using only three control-flow primitives: sequence,
selection (\texttt{if}), and a single repetition construct.  This result, together
with Dijkstra's influential argument against unstructured jumps \cite{dijkstra1968goto},
established the primacy of structured loops in procedural programming.
A practical consequence is that the four loop forms commonly found in procedural
languages --- \textit{for}, \textit{while}, \textit{do-while}, and
\textit{repeat-until} --- are mutually interconvertible: any loop can be rewritten using any other form, at the
cost of at most a guard condition or an unrolled first iteration.

This is not merely an academic curiosity.  The Romanian high-school computer science
curriculum explicitly teaches all four loop forms and expects students to
identify and apply these transformations in examination problems.
Providing an automated, interactive converter therefore has direct pedagogical value.

\subsection{Semantics of the Four Loop Forms}

Let $B$ denote the loop body (a finite sequence of statements), $C$ a Boolean
condition, $i$ a counter variable, $a$ the initial value, $b$ the bound, and $s > 0$
the step size.  The four loops have the following denotational semantics:

\begin{itemize}
  \item \textbf{for}: execute $B$ for $i = a,\, a{+}s,\, \ldots$ while $i \leq b$;
        zero or more iterations.
  \item \textbf{while}: while $C$, execute $B$; zero or more iterations.
  \item \textbf{do-while}: execute $B$, then repeat while $C$; at least one
        iteration.
  \item \textbf{repeat-until}: execute $B$, then repeat until $C$
        (i.e.\ while $\neg C$); at least one iteration.
\end{itemize}

The key structural distinction is between \textit{pre-test} loops (\textit{while},
\textit{for}) that may execute zero times, and \textit{post-test} loops
(\textit{do-while}, \textit{repeat-until}) that always execute at least once.
Crossing the boundary between the two families requires a guard.

\subsection{Transformation Rules}

The following semantics-preserving rewrite rules hold unconditionally:

\paragraph{while $\to$ do-while (and back).}
\[
  \textbf{while } C \;\textbf{do}\; B
  \;\;\equiv\;\;
  \textbf{if } C \;\textbf{then}\; (\textbf{do}\; B \;\textbf{while}\; C)
\]
\[
  \textbf{do}\; B \;\textbf{while}\; C
  \;\;\equiv\;\;
  B\,;\;\textbf{while } C \;\textbf{do}\; B
\]

\paragraph{for $\to$ while.}
\[
  \textbf{for}\; i \leftarrow a,\, b,\, s \;\textbf{do}\; B
  \;\;\equiv\;\;
  i \leftarrow a\,;\;\textbf{while}\; i \leq b \;\textbf{do}\; (B\,;\;i \leftarrow i + s)
\]

\paragraph{repeat-until $\to$ while (and back).}
\[
  \textbf{repeat}\; B \;\textbf{until}\; C
  \;\;\equiv\;\;
  B\,;\;\textbf{while}\; \neg C \;\textbf{do}\; B
\]

By composing these four base rules, all twelve directed pairwise conversions among
the four loop forms can be derived.  Pseudo++'s \texttt{equivalence} module
implements each conversion as a text-level rewrite over the CST, preserving the
original source indentation.

% ─────────────────────────────────────────────────────────────────────────────
\section{WebAssembly}
\label{sec:webassembly}

\subsection{Overview}

WebAssembly (Wasm) \cite{haas2017wasm,wasm-spec} is a binary instruction
format for a stack-based virtual machine, standardised by the W3C in 2019.  Its design goals are
portability, safety, and near-native execution speed in the browser.  A WebAssembly
module is a self-contained binary (\texttt{.wasm}) that exports and imports named
functions; the host (typically a JavaScript runtime) provides imports and calls
exports.

Key properties of the Wasm execution model:

\begin{itemize}
  \item \textbf{Linear memory}: a contiguous, byte-addressable memory region whose
        size is a multiple of 64~KiB.  The C heap and stack both reside in linear
        memory; pointer arithmetic is valid within this region.

  \item \textbf{Type safety}: all instructions are typed; the validator rejects
        ill-typed modules before execution.

  \item \textbf{Sandboxing}: a Wasm module cannot access host memory outside its
        own linear memory or perform arbitrary I/O; all interaction with the outside
        world happens through explicitly imported functions.

  \item \textbf{Determinism}: aside from NaN bit-patterns, Wasm execution is fully
        deterministic.
\end{itemize}

\subsection{Emscripten}

Emscripten \cite{emscripten} is an LLVM-based compiler toolchain that compiles C
and C++ to WebAssembly together with a JavaScript glue layer.  Given a C source
file, Emscripten compiles it to LLVM IR via \texttt{clang}, lowers it to Wasm
via \texttt{wasm-ld} and \texttt{wasm-opt}, and generates a JavaScript companion
(\texttt{.js}) that instantiates the module, provides standard C library shims,
and exposes exported C functions to JavaScript callers.

Functions to be exported must be annotated with \texttt{EMSCRIPTEN\_KEEPALIVE}
to prevent dead-code elimination, or listed via the
\texttt{-sEXPORTED\_FUNCTIONS} linker flag.


% ─────────────────────────────────────────────────────────────────────────────
\section{Statement-Level Step-Through Debugging}
\label{sec:debugging}

\subsection{Interactive Debugging in Educational Contexts}

Step-through debugging allows a student to observe a program's behaviour one
statement at a time, inspecting variable values and control flow at each point.
This is particularly valuable when learning programming: rather than reading static
code, the student sees the program \textit{in motion} --- which variable changes,
which branch is taken, how a loop counter evolves \cite{wing2006ct}.
More advanced tools support stepping \emph{backwards} through execution history
(omniscient or reversible debugging \cite{lewis2003omniscient,reversible-debugging});
Pseudo++ provides forward-only stepping.

Pseudo++'s debugger exposes two user-facing operations:

\begin{itemize}
  \item \textbf{Step Into} (one statement at a time): executes the next statement,
        highlights the corresponding source line in the editor, and updates the
        variable inspection panel.
  \item \textbf{Continue} (run to completion): resumes fast execution without
        step-by-step visibility, stopping only when the program finishes, encounters
        a runtime error, or needs user input.
\end{itemize}

